\begin{figure}[ht]
		\centering
		\begin{tikzpicture}[>=Stealth, line width=0.8pt]
			
			\def\R{2.0}
			\def\arr{0.55}
			\def\dx{1.4}    % separazione tra C_1 e C_2 nei corpi liberi
			
			% ============================================================
			% PANNELLO (a) — dati del problema
			% ============================================================
			\begin{scope}[xshift=0cm]
				
				% semiarchi
				\draw[line width=1.5pt] (-\R,0)
				arc[start angle=180, end angle=90, radius=\R];
				\draw[line width=1.5pt] (0,\R)
				arc[start angle=90, end angle=0, radius=\R];
				
				% cerniera in A
				\draw[thick] (-\R,0) -- (-\R-0.22,-0.38)
				-- (-\R+0.22,-0.38) -- cycle;
				\draw[thick] (-\R-0.32,-0.38) -- (-\R+0.32,-0.38);
				\foreach \x in {-0.24,-0.12,0,0.12,0.24}{
					\draw[thin] (-\R+\x,-0.38) -- (-\R+\x-0.10,-0.52);
				}
				% cerniera in B
				\draw[thick] (\R,0) -- (\R-0.22,-0.38)
				-- (\R+0.22,-0.38) -- cycle;
				\draw[thick] (\R-0.32,-0.38) -- (\R+0.32,-0.38);
				\foreach \x in {-0.24,-0.12,0,0.12,0.24}{
					\draw[thin] (\R+\x,-0.38) -- (\R+\x-0.10,-0.52);
				}
				
				% perni A, B, D
				\node[circle, fill=white, draw=black, inner sep=0pt,
				minimum size=5pt, line width=0.8pt] at (-\R,0) {};
				\node[circle, fill=white, draw=black, inner sep=0pt,
				minimum size=5pt, line width=0.8pt] at (\R,0) {};
				\node[circle, fill=white, draw=black, inner sep=0pt,
				minimum size=5pt, line width=0.8pt] at (0,\R) {};
				
				% etichette punti
				\node[left, font=\scriptsize] at (-\R,+0.2) {$A$};
				\node[right, font=\scriptsize] at (\R,+0.2) {$B$};
				\node[font=\scriptsize] at (-0.,\R+0.30) {$D$};
				
				% etichette corpi (esterne)
				\node[font=\small] at ({(\R+0.35)*cos(150)},{(\R+0.35)*sin(150)})	{$\mathcal{C}_1$};
				\node[font=\small] at ({(\R+0.35)*cos(35)},{(\R+0.35)*sin(35)})		{$\mathcal{C}_2$};
				
				% punto P_1 (centro asse C_1, a 135°)
				\node[circle, fill=black, inner sep=0pt, minimum size=2pt]
				at ({\R*cos(135)},{\R*sin(135)}) {};
				\node[font=\scriptsize] at ({(\R-0.30)*cos(135)},{(\R-0.30)*sin(135)+0.07}) {$P_1$};
				% coppia mu antioraria in P_1
				\draw[->, thick] ({\R*cos(135)+0.30},{\R*sin(135)+0.05}) arc[start angle=10, end angle=310, radius=0.30];
				\node[font=\scriptsize]	at ({\R*cos(135)+0.55},{\R*sin(135)+0.10}) {$\mu$};
				
				% punto P_2 (centro asse C_2, a 45°)
				\node[circle, fill=black, inner sep=0pt, minimum size=2pt]
				at ({\R*cos(45)},{\R*sin(45)}) {};
				\node[font=\scriptsize]	at ({(\R-0.30)*cos(45)},{(\R-0.30)*sin(45)}) {$P_2$};
				% forza F verticale verso il basso in P_2
				\draw[->, thick] ({\R*cos(45)},{\R*sin(45)+\arr}) -- ({\R*cos(45)},{\R*sin(45)});
				\node[font=\scriptsize]	at ({\R*cos(45)+0.25},{\R*sin(45)+\arr-0.10}) {$F$};
				
				
				
				% etichetta pannello
				\node[font=\scriptsize] at (0,-0.50) {(a)};
				
			\end{scope}
			
			% ============================================================
			% PANNELLO (b) — schemi dei corpi liberi
			% ============================================================
			\begin{scope}[xshift=7.0cm]
				
				% ----- C_1 libero (decentrato a sinistra) -----
				\begin{scope}[xshift=-\dx cm]
					
					% semiarco C_1
					\draw[line width=1.5pt] (-\R,0)
					arc[start angle=180, end angle=90, radius=\R];
					
					% cerniera in A (grigia, vincolo rotto)
					\draw[thick, gray!50] (-\R,0) -- (-\R-0.22,-0.38)
					-- (-\R+0.22,-0.38) -- cycle;
					\draw[thick, gray!50] (-\R-0.32,-0.38) -- (-\R+0.32,-0.38);
					\foreach \x in {-0.24,-0.12,0,0.12,0.24}{
						\draw[thin, gray!50] (-\R+\x,-0.38) -- (-\R+\x-0.10,-0.52);
					}
					
					% perni A, D
					\node[circle, fill=white, draw=black, inner sep=0pt,
					minimum size=5pt, line width=0.8pt] at (-\R,0) {};
					\node[circle, fill=white, draw=black, inner sep=0pt,
					minimum size=5pt, line width=0.8pt] at (0,\R) {};
					
					% etichette
					\node[left, font=\scriptsize] at (-\R,-0.) {$A$};
					\node[below=2pt, font=\scriptsize] at (0,\R+0.05) {$D$};
					\node[font=\small] at ({(\R+0.35)*cos(150)},{(\R+0.35)*sin(150)})	{$\mathcal{C}_1$};
					
					
					
					% reazioni in A
					\draw[->, thick] (-\R,0.0) -- (-\R+\arr,0.0)
					node[above, font=\scriptsize] {$X_A$};
					\draw[->, thick] (-\R-0.05,0) -- (-\R-0.05,\arr)
					node[left, yshift=-3pt, font=\scriptsize] {$Y_A$};
					
					% reazioni interne in D (su C_1): X_D verso destra, Y_D verso l'alto
					\draw[->, thick] (0,\R) -- (\arr,\R)
					node[below, font=\scriptsize] {$X_D$};
					\draw[->, thick] (-0.05,\R) -- (-0.05,\R+\arr)
					node[left, yshift=-3pt, font=\scriptsize] {$Y_D$};
					
					% coppia mu in P_1
					\node[circle, fill=black, inner sep=0pt, minimum size=2pt]
					at ({\R*cos(135)},{\R*sin(135)}) {};
					\draw[->, thick]
					({\R*cos(135)+0.25},{\R*sin(135)+0.05})
					arc[start angle=10, end angle=310, radius=0.25];
					\node[font=\scriptsize]	at ({\R*cos(135)+0.50},{\R*sin(135)+0.05}) {$\mu$};
					\node[font=\scriptsize] at ({(\R-0.30)*cos(135)},{(\R-0.30)*sin(135)+0.10}) {$P_1$};
					
				\end{scope}
				
				% ----- C_2 libero (decentrato a destra) -----
				\begin{scope}[xshift=\dx-0.10 cm]
					
					% semiarco C_2
					\draw[line width=1.5pt] (0,\R)
					arc[start angle=90, end angle=0, radius=\R];
					
					% cerniera in B (grigia)
					\draw[thick, gray!50] (\R,0) -- (\R-0.22,-0.38)
					-- (\R+0.22,-0.38) -- cycle;
					\draw[thick, gray!50] (\R-0.32,-0.38) -- (\R+0.32,-0.38);
					\foreach \x in {-0.24,-0.12,0,0.12,0.24}{
						\draw[thin, gray!50] (\R+\x,-0.38) -- (\R+\x-0.10,-0.52);
					}
					
					% perni B, D
					\node[circle, fill=white, draw=black, inner sep=0pt,
					minimum size=5pt, line width=0.8pt] at (\R,0) {};
					\node[circle, fill=white, draw=black, inner sep=0pt,
					minimum size=5pt, line width=0.8pt] at (0,\R) {};
					
					% etichette
					\node[left, font=\scriptsize] at (\R,+0.) {$B$};
					\node[above, font=\scriptsize] at (0,\R) {$D$};
					\node[font=\small] at ({(\R+0.35)*cos(35)},{(\R+0.35)*sin(35)})	{$\mathcal{C}_2$};
					
					
					% reazioni in B
					\draw[->, thick] (\R,0.0) -- (\R+\arr,0.0)
					node[above, font=\scriptsize] {$X_B$};
					\draw[->, thick] (\R+0.05,0) -- (\R+0.05,\arr)
					node[left, yshift=-3pt, font=\scriptsize] {$Y_B$};
					
					% reazioni interne in D (su C_2): opposte (verso sinistra e basso)
					\draw[->, thick] (0,\R) -- (-\arr,\R)
					node[above, font=\scriptsize] {$X_D$};
					\draw[->, thick] (0.05,\R) -- (0.05,\R-\arr)
					node[right, yshift=3pt, font=\scriptsize] {$Y_D$};
					
					% forza F in P_2
					\node[circle, fill=black, inner sep=0pt, minimum size=2pt]
					at ({\R*cos(45)},{\R*sin(45)}) {};
					\draw[->, thick] ({\R*cos(45)},{\R*sin(45)+\arr}) -- ({\R*cos(45)},{\R*sin(45)});
					\node[font=\scriptsize]	at ({\R*cos(45)+0.25},{\R*sin(45)+\arr-0.10}) {$F$};
					\node[font=\scriptsize]	at ({(\R-0.30)*cos(45)},{(\R-0.30)*sin(45)}) {$P_2$};
					
				\end{scope}
				
				% etichetta pannello
				\node[font=\scriptsize] at (-0.60,-0.50) {(b)};
				
			\end{scope}
			
		\end{tikzpicture}
		\caption{Arco a tre cerniere
			(\ESEref{ex:arco_statico}):
			\emph{(a)}~dati del problema statico ---
			carico verticale $F$ in $P_2$ (punto medio
			dell'asse di $\mathcal{C}_2$) e coppia $\mu$
			antioraria in $P_1$ (punto medio dell'asse
			di $\mathcal{C}_1$); le frecce indicano il
			verso positivo dei carichi assegnati.
			\emph{(b)}~schemi dei corpi liberi di
			$\mathcal{C}_1$ (sinistra) e $\mathcal{C}_2$
			(destra): reazioni esterne $X_A, Y_A$ in $A$
			e $X_B, Y_B$ in $B$; reazioni interne $X_D, Y_D$
			in $D$, uguali e opposte sui due corpi.}
		\label{fig:arco_statico}
	\end{figure}
